% =====================================================================
%  12  プロピレン精製 — 膜ユニット（2段：上=概念図/プロファイル、下=表/膜選定）
%  source_memo: 07章＋top1_trial201.txt[Membrane]＋Hua et al. ACS AMI 2024,16,15273.
%   ZIF-8細孔4.0-4.2Å。多結晶ZIF-8膜 SF~300/52GPU(大面積化・高コスト難)。ポリマー膜=トレードオフ・可塑化。
%   ZIF-8 混合マトリックス膜(添加剤併用in-situ形成): 無添加57.7/22.5・最高222.5/10.1・頑健量産版74.9/26.6GPU。採用α90/40=楽観側。
%  方針：図で一目＋口頭。膜=緑SepGreen。丸括弧禁止。余白を作らず最下部まで使う。
% =====================================================================
\definecolor{SepGreen}{HTML}{2E9E5B}
\definecolor{DipeOrange}{HTML}{E8821E}
\begin{frame}[t]

  \vspace{0.1em}

  % ===== 上段：膜分離器の概念図（左）／膜面積プロファイル（右）=====
  \begin{columns}[T,onlytextwidth]
    \begin{column}{0.45\textwidth}
      \raggedright
      {\footnotesize\bfseries\color{SepGreen!45!black} 膜分離器}\par
      \vspace{0.4em}
      \resizebox{\linewidth}{!}{%
      \begin{tikzpicture}[
        x=1cm, y=1cm, >=Stealth, font=\Large,
        c6/.style={circle, fill=SepGreen, draw=black, line width=0.4pt, minimum size=4mm, inner sep=0},
        c8/.style={circle, fill=DipeOrange, draw=black, line width=0.4pt, minimum size=6mm, inner sep=0},
        io/.style={align=center, font=\Large},
        pass/.style={-{Stealth[length=2.2mm]}, SepGreen!70!black, line width=1.1pt},
      ]
        % 膜（細孔つき）
        \fill[SepGreen!28] (2.4,0.2) rectangle (2.9,3.6);
        \draw[SepGreen, line width=1pt] (2.4,0.2) rectangle (2.9,3.6);
        \foreach \y in {0.7,1.4,2.1,2.8,3.3}{\fill[white] (2.4,\y-0.11) rectangle (2.9,\y+0.11);}
        \node[SepGreen!45!black, font=\Large\bfseries, anchor=south] at (2.65,3.7) {ZIF-8 膜};
        % 供給側（左）の分子：C3H8(橙)は残る、C3H6(緑)は透過
        \node[c8] at (1.0,2.9){}; \node[c6] at (1.75,3.1){}; \node[c8] at (1.3,2.0){};
        \node[c6] at (2.0,1.7){}; \node[c8] at (1.05,1.05){}; \node[c6] at (1.8,0.8){};
        % C3H6 が細孔を透過
        \draw[pass] (2.1,2.5)--(3.45,2.5); \draw[pass] (2.2,1.7)--(3.45,1.7); \draw[pass] (2.1,0.85)--(3.45,0.85);
        \node[c6] at (3.65,2.5){}; \node[c6] at (3.75,1.7){}; \node[c6] at (3.6,0.85){};
        % 供給（左、P_H 併記で重なり回避）
        \node[io, anchor=east] (f) at (-0.05,3.0) {供給 $\mathrm{C_3H_6}/\mathrm{C_3H_8}$\\[1pt]{\normalsize\bfseries\color{HeaderBlue}$P_H$ 8.4 bar}};
        \draw[-{Stealth[length=2.6mm]}, HeaderBlue, line width=1.4pt] (f.east) -- (0.7,3.05);
        % 非透過（下へ・矢印の下にラベル）
        \draw[-{Stealth[length=2.6mm]}, DipeOrange!80!black, line width=1.3pt] (1.5,0.0) -- (1.5,-0.75);
        \node[io, anchor=north, text=DipeOrange!70!black] at (1.5,-0.85) {非透過 $\mathrm{C_3H_8}$ → 循環};
        % 透過（右、98.6% 大・P_L 併記）
        \draw[pass, line width=1.5pt] (4.0,1.7)--(4.9,1.7);
        \node[io, anchor=west, text=SepGreen!45!black] at (4.95,2.15) {透過 $\mathrm{C_3H_6}$};
        \node[io, anchor=west, font=\LARGE\bfseries, text=SepGreen!45!black] at (4.95,1.45) {98.6\%};
        \node[io, anchor=west, font=\normalsize\bfseries, text=HeaderBlue] at (4.95,0.78) {$P_L$ 1 atm};
      \end{tikzpicture}}
    \end{column}

    \begin{column}{0.55\textwidth}
      \centering
      \includegraphics[width=\linewidth]{membrane_profile}\par
      \vspace{-0.82em}
      {\scriptsize モジュール内での $\mathrm{C_3H_6}$ 組成の変化}
    \end{column}
  \end{columns}

  \vspace{0.75em}

  % ===== 下段：設計変数表（左）／膜の選定 枠囲み（右）==========
  \begin{columns}[T,onlytextwidth]
    \begin{column}{0.37\textwidth}
      {\scriptsize\bfseries\color{HeaderBlue} 設計変数と性能}\par
      \vspace{0.12em}
      {\scriptsize\setlength{\tabcolsep}{5pt}\renewcommand{\arraystretch}{0.95}%
       \begin{tabular}{@{}l r@{}}
         \toprule
         供給圧 $P_H$ & $8.4$ bar \\
         膜面積 $A_{\mathrm{mem}}$ & $1.28\times10^{5}$ m$^2$ \\
         \midrule
         選択性 $\alpha$ & $90$ \\
         透過度 $Q_A$ & $40$ GPU \\
         透過圧 $P_L$ & $1$ atm \\
         \midrule
         モジュール数 & $256$ 本 \\
         ステージカット & $37$ \% \\
         透過純度 $\mathrm{C_3H_6}$ & $98.6$ mol\% \\
         \bottomrule
       \end{tabular}}\par
      \vspace{0.12em}
      {\tiny\color{black!55} 上2行が最適化変数\par
       \vspace{0.10em}
       GPU：気体透過率の単位\par
       \smash{$1\,\mathrm{GPU}=3.35\times10^{-10}\ \mathrm{mol\,m^{-2}s^{-1}Pa^{-1}}$}\par
       \vspace{0.6em}
       {\color{black!50} 出典　Hua et al., ACS AMI, 16, 15273, 2024}}
    \end{column}

    \begin{column}{0.63\textwidth}
      {\setlength{\fboxsep}{6pt}\setlength{\fboxrule}{1.2pt}%
       \fcolorbox{SepGreen!55!black}{SepGreen!6}{%
        \begin{minipage}[t][4.0cm]{0.93\linewidth}
         {\scriptsize\bfseries\color{SepGreen!40!black} 膜の選定　ZIF-8 を高分子に分散した混合マトリックス膜}
         \vfill
         {\scriptsize\setlength{\tabcolsep}{4pt}\renewcommand{\arraystretch}{1.1}%
          \begin{tabular}{@{}l c c l@{}}
            \toprule
            膜種 & $\alpha$ & 透過 GPU & 特徴 \\
            \midrule
            多結晶 ZIF-8 & $\sim\!300$ & $52$ & 大面積化・高コスト \\
            高分子膜 & 低 & --- & 可塑化・低性能 \\
            {\bfseries\color{SepGreen!40!black} 採用 混合膜} & {\bfseries 90} & {\bfseries 40} & 量産可・高選択 \\
            \bottomrule
          \end{tabular}}
         \vfill
         {\scriptsize\raggedright 高選択の鍵\par \textbf{小さい非晶質 ZIF-8}・添加剤併用の in-situ 形成で量産\par}
         \vfill
         {\scriptsize\raggedright {\color{SpRed}\textbf{!}}\ 実験室レベル　大面積・連続化で低下＝採用値は楽観側\par}
         \vfill
         {\scriptsize\raggedright 実機の劣化要因：可塑化・競争吸着・濃度分極 等\par}
        \end{minipage}}}
    \end{column}
  \end{columns}

\end{frame}
